\documentclass[12pt]{article}
\usepackage{amsmath}
\usepackage{amssymb}
\usepackage{amsthm}
\usepackage{geometry}
\geometry{margin=1.4in}
\usepackage{hyperref}
\usepackage{setspace}
\setstretch{1.4}

\newtheorem{theorem}{Theorem}
\newtheorem{conjecture}[theorem]{Conjecture}
\theoremstyle{definition}
\newtheorem{definition}[theorem]{Definition}
\theoremstyle{remark}
\newtheorem{remark}[theorem]{Remark}

\title{The Missing Frobenius}
\author{Diego Rinc\'on\\
\small Independent}
\date{}

\begin{document}

\maketitle

\begin{abstract}
Over a finite field $\mathbb{F}_q$,
there is a canonical symmetry: the Frobenius $x \mapsto x^q$.
It is the reason the Weil conjectures are provable.
Over $\mathbb{Q}$, there is no global Frobenius.
This is the reason the Riemann Hypothesis is not provable by the same method.
The missing Frobenius is not a conjecture. It is an absence.
What is absent, and why, and what its presence would mean.
\end{abstract}

\section{What the Frobenius Is}

Over the finite field $\mathbb{F}_q$ with $q = p^n$ elements,
the Frobenius morphism is:
$$\varphi_q : x \longmapsto x^q.$$
It fixes exactly the elements of $\mathbb{F}_q$
and permutes everything else.
It is a canonical element of $\mathrm{Gal}(\overline{\mathbb{F}_q}/\mathbb{F}_q)$:
the generator, the clock of the field.

For a variety $X/\mathbb{F}_q$,
the Frobenius acts on the \'etale cohomology $H^*_{\rm\acute{e}t}(X_{\bar{\mathbb{F}}_q}, \mathbb{Q}_\ell)$.
Its eigenvalues are the arithmetic of $X$:
the number of points over $\mathbb{F}_{q^n}$
is determined by the trace of $\varphi_q^n$ on cohomology.

This is the Lefschetz trace formula, and it is exact.
The Frobenius encodes everything.

\begin{theorem}[Weil, Deligne {\cite{deligne,weil}}]
For a smooth projective variety $X/\mathbb{F}_q$,
the eigenvalues of $\varphi_q$ on $H^i_{\rm\acute{e}t}(X, \mathbb{Q}_\ell)$
are algebraic integers of absolute value $q^{i/2}$.
In particular, the zeros of the zeta function of $X$
lie on the lines $\mathrm{Re}(s) = i/2$
for $i = 0, 1, \ldots, 2\dim X$.
\end{theorem}

This is the Riemann Hypothesis over finite fields.
It is proved.
It is proved because the Frobenius exists.

\section{Why There Is No Global Frobenius}

Over $\mathbb{Q}$, there is no analogue.

At each prime $p$, there is a local Frobenius:
$\varphi_p \in \mathrm{Gal}(\bar{\mathbb{Q}}_p/\mathbb{Q}_p)$,
well-defined up to conjugacy.
The local Frobenii are the building blocks of the Galois group $\mathrm{Gal}(\bar{\mathbb{Q}}/\mathbb{Q})$.
But there is no single element of the global Galois group
that acts as the Frobenius everywhere simultaneously.

The reason is structural.
Over $\mathbb{F}_q$, the Galois group is $\hat{\mathbb{Z}}$,
generated by one element ($\varphi_q$).
Over $\mathbb{Q}$, the Galois group $\mathrm{Gal}(\bar{\mathbb{Q}}/\mathbb{Q})$
is profinite, non-abelian, and vastly more complex.
It has no preferred generator.
It has local Frobenii at each prime,
but they do not cohere into a global one.

\begin{remark}
This is not a theorem that says ``the global Frobenius does not exist.''
It is an observation about the structure of $\mathrm{Gal}(\bar{\mathbb{Q}}/\mathbb{Q})$:
it has no element with the properties the Frobenius would need to have.
The absence is structural, not accidental.
\end{remark}

\section{What the Frobenius Would Do}

If a global Frobenius $\varphi_{\mathbb{Q}}$ existed,
acting on the motivic cohomology $H^*_{\rm mot}(\mathrm{Spec}\,\mathbb{Q})$
with the right eigenvalue properties,
the following would follow:

\begin{enumerate}
\item \textbf{Riemann Hypothesis.}
The zeros of $\zeta(s)$ would be the eigenvalues of $\varphi_{\mathbb{Q}}$
on $H^1_{\rm mot}(\mathrm{Spec}\,\mathbb{Z})$.
By the Weil-type bound, they would have absolute value $q^{1/2}$
for the appropriate $q$ --- placing them on $\mathrm{Re}(s) = 1/2$.

\item \textbf{Birch--Swinnerton-Dyer.}
The rank of $E(\mathbb{Q})$ would be the multiplicity of the eigenvalue $1$
of $\varphi_{\mathbb{Q}}$ on $H^1_{\rm mot}(E)$.
The order of vanishing of $L(E,s)$ at $s=1$ would be this multiplicity.

\item \textbf{Hodge conjecture (over $\mathbb{Q}$).}
The global Frobenius would distinguish algebraic classes
from transcendental ones:
an algebraic cycle has Frobenius eigenvalue a root of unity;
a transcendental class does not.
The $(p,p)$ condition would follow from Frobenius acting as
multiplication by $q^p$.
\end{enumerate}

Three walls.
One missing symmetry.

\section{Candidates for the Missing Frobenius}

The absence has been noticed.
Several structures have been proposed as its replacement or incarnation.

\textbf{The Fargues--Fontaine curve.}
Fargues and Fontaine constructed a curve $X_{\rm FF}$
such that $p$-adic representations of $\mathrm{Gal}(\bar{\mathbb{Q}}_p/\mathbb{Q}_p)$
correspond to vector bundles on $X_{\rm FF}$.
The Frobenius on $X_{\rm FF}$ is the $p$-adic Frobenius.
This works locally. The global version is not yet constructed.

\textbf{The $\Lambda$-ring structure.}
The ring $\mathbb{Z}$ carries a $\Lambda$-ring structure
(Adams operations $\psi^p$ for each prime $p$)
which behaves like a family of Frobenii.
This is the basis of the theory of $\Lambda$-schemes
and related to the ``field with one element'' $\mathbb{F}_1$.
A geometry over $\mathbb{F}_1$ would have a single Frobenius
reducing to $\varphi_p$ at each prime.

\textbf{Connes' scaling action.}
In the Bost--Connes system \cite{bostconnes},
the time evolution $\sigma_t(f) = f \cdot |\cdot|^{it}$
plays the role of the Frobenius.
The KMS states at inverse temperature $\beta > 1$
are indexed by $\hat{\mathbb{Z}}^* \cong \mathrm{Gal}(\mathbb{Q}^{\rm ab}/\mathbb{Q})$.
At $\beta = 1$ (the phase transition), the symmetry breaks.
This is not quite a Frobenius, but it is the closest
operator-algebraic analogue.

\textbf{The absolute Galois group.}
The deepest version: $\mathrm{Gal}(\bar{\mathbb{Q}}/\mathbb{Q})$ itself
might act on the motivic cohomology of $\mathrm{Spec}\,\mathbb{Z}$
in a way that recovers the Frobenius at each prime.
This is what the Langlands program would give,
if the global Langlands correspondence were fully established.
The global Frobenius would not be a single element
but a representation of the whole group.

\section{The Grammar of the Absence}

The grammar paper \cite{grammar} identifies the local-to-global passage
as the wall.
The missing Frobenius is the reason for the wall.

Over $\mathbb{F}_q$: the local and global Frobenii coincide.
The grammar is complete.
The Riemann Hypothesis is provable.

Over $\mathbb{Q}$: the local Frobenii exist but do not cohere globally.
The grammar is incomplete.
The Riemann Hypothesis is not provable by the same method.

What is needed is not a proof within the current grammar.
What is needed is a global Frobenius ---
or something that does what the global Frobenius would do:
a canonical symmetry of the arithmetic of $\mathbb{Q}$
that acts on the motivic cohomology with the right eigenvalue properties.

When that symmetry is found ---
if it is found ---
the proof will be the coboundary it generates.
The wall will not be overcome.
It will dissolve.
The grammar will have been extended to include what was absent.

\begin{thebibliography}{9}

\bibitem{deligne}
Deligne, P. (1974).
La conjecture de Weil, I.
\textit{Publications Math\'ematiques de l'IH\'ES}, 43, 273--307.

\bibitem{weil}
Rinc\'on, D. (2026).
The Weil Conjectures: Proof of Concept for RH. Preprint.

\bibitem{bostconnes}
Rinc\'on, D. (2026).
Bost--Connes: Zeta as Partition Function. Preprint.

\bibitem{grammar}
Rinc\'on, D. (2026).
Mathematics as Cohomological Grammar. Preprint.

\bibitem{one_wall}
Rinc\'on, D. (2026).
One Wall: The Local-to-Global Conjecture. Preprint.

\bibitem{langlands}
Rinc\'on, D. (2026).
The Langlands Thread. Preprint.

\bibitem{perfectoid}
Rinc\'on, D. (2026).
Perfectoid Spaces and the Local Frobenius. Preprint.

\end{thebibliography}

\end{document}
